\section{Introduction}

Retrieval infrastructure is becoming operational memory for agentic AI systems. Agents repeatedly query, update, and reuse external stores while operating under latency budgets and context-window costs, so a retrieval engine that looks strong on a static relevance benchmark may be unsuitable when inserts must become visible quickly, tail latency must remain bounded, or retrieved passages dominate downstream model cost. This paper studies the narrower question that an operator must answer before deploying an agentic retrieval stack: can a commodity-hardware benchmark produce stable deployment decisions for retrieval infrastructure?

Existing evidence is fragmented across two evaluation styles. Information-retrieval benchmarks such as BEIR \citep{thakur2021beir} and HotpotQA \citep{yang2018hotpotqa} provide reusable relevance judgments, but they do not specify engine conformance, update visibility, or tail-latency constraints. Systems benchmarks often report throughput and latency, but they may not connect those measurements to retrieval quality, freshness after writes, or the token cost induced by retrieved contexts. This gap makes it easy to overfit a deployment decision to a single metric: a system can be fastest, cheapest, or highest quality under different objectives.

\sys addresses this gap with a reproducible, single-node benchmark for agentic retrieval infrastructure. The benchmark uses a fixed adapter contract, conformance tests before reporting, three workload scenarios, two local embedding tiers, and a budget ladder that measures whether decisions are stable as runs become longer. The central design choice is to treat benchmarking as a deployment decision audit rather than only a leaderboard: for each workload, \sys identifies the minimum viable engine and embedding under a strict p99 latency threshold, separately reports unconstrained-cost and quality-first alternatives, and exposes the margins that make a decision robust or fragile.

The archived evaluation gives three concrete findings. First, under a 200 ms worst-case p99 threshold, FAISS CPU with bge-small is the minimum viable deployment for single-hop, streaming-memory, and multi-hop evidence retrieval. The decision-margin audit shows that this result should be read locally: S1 and S2 are cost/quality ties resolved by tail latency, whereas S3 has lower cost and higher evidence coverage than the next strict-cost candidate. Targeted paired audits check the most likely counter-stories: a larger same-machine S2 run finds no hidden Qdrant quality or freshness win, and a 5{,}000-query S3 audit demotes the pgvector quality-first row to objective-sensitivity evidence rather than a substantive quality advantage. Second, changing the objective can change the selected engine or embedding tier, so the benchmark should report strict-latency, unconstrained-cost, and quality-first decisions separately. Third, decision stability is workload-dependent: B0-to-B2 top-1 agreement is absent for S1 and S2 but present for S3, even though S3 has low full-rank Spearman correlation. These results support a cautious conclusion: commodity-hardware benchmarks can produce useful deployment evidence, but only when the paper reports conformance, budgets, decision margins, and objective sensitivity explicitly.

This paper contributes: (i) a conformance-gated benchmark protocol for agentic retrieval infrastructure; (ii) three portable workloads covering static retrieval, streaming memory, and bounded multi-hop evidence retrieval; (iii) a deployment-decision analysis that separates strict-latency, cost, quality, decision-margin, and paired-audit evidence; and (iv) an archived result bundle with confidence intervals, behavior-card support tables, metadata, and reproducibility commands.
